
\section{The certificate-only principle}\label{sec:principle}

We state the rule formally and record what it forces.

\begin{definition}[Certificate]\label{def:cert}
Let $\phi$ be a feature of $\Gstar$---a proposition such as ``$i$ and $j$ are
adjacent'' or ``$\Gstar$ contains the arrowhead $i \to j$''. A
\emph{certifying null} for $\phi$ is a hypothesis $H_\phi$ with the property
that $P \in H_\phi$ whenever $\phi$ is false. A \emph{certificate for $\phi$
at level $\alpha$} is the event that an e-process for $H_\phi$ exceeds
$1/\alpha$.
\end{definition}

The implication in Definition~\ref{def:cert} runs one way only, and that is
the whole point. $H_\phi$ must be true whenever $\phi$ is false; it may also
be true in some cases where $\phi$ is true, and then the certificate simply
never fires. This asymmetry is what buys validity without faithfulness, and
what costs power.

\begin{proposition}[Soundness of a single certificate]\label{prop:single}
If $(E_t)$ is an e-process for a certifying null $H_\phi$, then
$\Prob(\exists t: \text{$\phi$ is certified at level $\alpha$ and $\phi$ is
false}) \le \alpha$.
\end{proposition}

\begin{proof}
If $\phi$ is false then $P \in H_\phi$ by Definition~\ref{def:cert}, so
$(E_t)$ is an e-process under $P$ and~\eqref{eq:ville} bounds the probability
that it ever exceeds $1/\alpha$ by $\alpha$. If $\phi$ is true the event in
question is empty.
\end{proof}

Proposition~\ref{prop:single} is elementary; we state it to make explicit that
the burden of the design is entirely in constructing $H_\phi$, not in the
inference. Three structural consequences follow, and they shape everything
below.

\paragraph{The algorithm cannot delete.} There is no certifying null for the
proposition ``$i$ and $j$ are non-adjacent''. Such a null would have to be
true whenever the pair \emph{is} adjacent, that is, it would have to be a
dependence hypothesis rejected by evidence of independence; but evidence of
independence is exactly what a hypothesis test does not provide. Equivalence
testing supplies the nearest available substitute---one can reject ``the
partial correlation exceeds $\delta$ in absolute value''---but the resulting
statement is ``the association is smaller than $\delta$'', which is a claim
about effect size, not about absence. We return to this in
Section~\ref{sec:negligible}. In the absence of such a null, an algorithm
obeying the rule must start from the empty graph.

\paragraph{The output is three-valued.} Each pair is in exactly one of three
states: \emph{certified adjacent} (with or without a certified orientation),
\emph{undecided}, or---if the extension of Section~\ref{sec:negligible} is
enabled---\emph{certified negligible at margin $\delta$}. The absence of an
edge in the output graph is not an assertion.

\paragraph{Power is descriptive, not guaranteed.} Nothing in
Proposition~\ref{prop:single} promises that a true edge is ever certified.
Whether it is depends on adjacency-faithfulness, on the sample size, and on
the primitive's power. This is a real limitation and we treat it as one:
Section~\ref{sec:experiments} reports certification rates alongside error
rates throughout, and reports the fraction of pairs left undecided rather than
suppressing it.

\section{The adjacency certificate}\label{sec:adjacency}

\subsection{The separability null}\label{sec:sepnull}

Fix a pair $(i,j)$ and consider
\begin{equation}\label{eq:sepnull}
H^{\mathrm{sep}}_{ij} \;:\quad
\exists\, S \in \Scal_k(i,j) \ \text{ such that } \ X_i \indep X_j \mid X_S .
\end{equation}

\begin{lemma}\label{lem:certifying}
Under Assumptions~\ref{ass:markov} and~\ref{ass:sep},
$H^{\mathrm{sep}}_{ij}$ is a certifying null for the adjacency of $i$ and $j$.
\end{lemma}

\begin{proof}
Suppose $i$ and $j$ are non-adjacent in $\Gstar$. By
Assumption~\ref{ass:sep} there is a set $S$ with $|S| \le k$ that d-separates
them, and $S \subseteq V \setminus \{i,j\}$, so $S \in \Scal_k(i,j)$. By
Assumption~\ref{ass:markov}, $X_i \indep X_j \mid X_S$ under $P$. Hence
$P \in H^{\mathrm{sep}}_{ij}$.
\end{proof}

Note what the lemma does \emph{not} require. It says nothing about adjacent
pairs. If $i$ and $j$ are adjacent, $P$ may or may not lie in
$H^{\mathrm{sep}}_{ij}$: an adjacent pair whose dependence is exactly
cancelled by a path through some $S$ satisfies the null too, and then the
certificate cannot fire. That is the failure of adjacency-faithfulness, and it
costs power, not validity.

\subsection{The minimum construction}\label{sec:minconstruction}

Let $(E^{(S)}_t)_{t \ge 0}$ be an e-process for the point null
$X_i \indep X_j \mid X_S$, for each $S \in \Scal_k(i,j)$, all built from the
same data stream. Define
\begin{equation}\label{eq:adjacency}
A_t(i,j) \;=\; \min_{S \in \Scal_k(i,j)} E^{(S)}_t(i,j) .
\end{equation}

\begin{proposition}\label{prop:adjacency}
$(A_t)_{t \ge 0}$ is an e-process for $H^{\mathrm{sep}}_{ij}$.
\end{proposition}

\begin{proof}
Let $P \in H^{\mathrm{sep}}_{ij}$. By~\eqref{eq:sepnull} there exists
$S^\star \in \Scal_k(i,j)$ with $X_i \indep X_j \mid X_{S^\star}$ under $P$,
so $P$ satisfies the point null for which $(E^{(S^\star)}_t)$ is an e-process.
By~\eqref{eq:adjacency}, $A_t \le E^{(S^\star)}_t$ pointwise for every $t$,
hence $A_\tau \le E^{(S^\star)}_\tau$ for every stopping time $\tau$, and
non-negativity gives
$\Exp_P[A_\tau] \le \Exp_P[E^{(S^\star)}_\tau] \le 1$.
\end{proof}

Combining Lemma~\ref{lem:certifying}, Proposition~\ref{prop:adjacency} and
Proposition~\ref{prop:single} gives the adjacency guarantee: the probability
that $A_t(i,j)$ ever exceeds $1/\alpha$ for a non-adjacent pair is at most
$\alpha$.

That a minimum of e-values is an e-value for a union null is standard
\citep{vovk2021evalues}, and the proof above is the standard one. What is
specific here is the identification of adjacency as the complement of such a
union, and the two consequences that follow.

\begin{remark}[No faithfulness]\label{rem:faithfulness}
The proof uses only that~\eqref{eq:sepnull} is true for non-adjacent pairs,
which follows from Assumptions~\ref{ass:markov} and~\ref{ass:sep}.
Faithfulness enters only through power: adjacency-faithfulness
\citep{ramsey2006adjacency} is what makes every $E^{(S)}$ grow for a genuinely
adjacent pair, and hence what makes the minimum grow. An instance where it
fails leaves pairs undecided rather than producing wrong edges. Stated
exactly, validity requires three things: the Markov condition,
Assumption~\ref{ass:sep}, and a valid per-set e-process. The third is not
free. Our default primitive (Section~\ref{sec:primitive}) is a linear-Gaussian
safe test, so that model is assumed \emph{there}; substituting a nonparametric
sequential test \citep{he2026sequential} removes it at a cost in power.
Proposition~\ref{prop:adjacency} is indifferent to which primitive is used,
and Section~\ref{sec:exp-primitive} measures the primitive's calibration under
six noise laws, including three outside its model.
\end{remark}

\begin{remark}[Multiplicity over pairs]\label{rem:multiplicity-pairs}
The construction tests one hypothesis per pair, not one per (pair,
conditioning set) triple. A union bound over the adjacency family therefore
runs over $\binom{d}{2}$ hypotheses rather than
$\binom{d}{2} \cdot |\Scal_k|$. At $d = 6$ and $k = 2$ we have
$|\Scal_k| = \binom{4}{0} + \binom{4}{1} + \binom{4}{2} = 11$, so the
per-hypothesis level is $11$ times larger than a triple-indexed budget would
allow. The saving grows with $k$: at $d = 10$, $k = 3$ it is $93$. This is not
a free lunch---the minimum in~\eqref{eq:adjacency} is itself conservative,
since it discards all but the least favourable conditioning set---but the
conservatism is of a different kind from a multiplicity correction, and it is
the kind that vanishes as the least favourable set becomes uninformative.
\end{remark}

\subsection{The per-set primitive}\label{sec:primitive}

Proposition~\ref{prop:adjacency} needs an e-process for each point null
$X_i \indep X_j \mid X_S$. For linear-Gaussian data we use an exact one,
obtained as a Bayes factor with right-Haar priors on the nuisance parameters.
Write the hypothesis in regression form,
\begin{equation}\label{eq:regform}
X_i \;=\; X_S^\top \alpha \;+\; \beta X_j \;+\; \varepsilon,
\qquad \varepsilon \sim N(0, \sigma^2),
\end{equation}
so that the null is $\beta = 0$ with nuisance parameters $(\alpha, \sigma)$
shared between the hypotheses. Place the right-Haar prior
$\pi(\alpha, \sigma) \propto \sigma^{-1}$ on the nuisance block under both
hypotheses, and $\beta \mid \sigma \sim N(0, \sigma^2 g)$ on the tested
coefficient under the alternative. Integrating both marginal likelihoods in
closed form (Appendix~\ref{app:safelinear}) gives, after $n$ observations,
\begin{equation}\label{eq:safelinear}
E_n \;=\; (1+G)^{-1/2}\,\bigl(1 - \kappa\,\hat r_n^2\bigr)^{-(n-p)/2},
\qquad
G \;=\; g\,S_{yy}, \quad \kappa \;=\; \frac{G}{1+G},
\end{equation}
where $p = |S| + 1$; $S_{xx}$, $S_{xy}$, $S_{yy}$ are the cross-products of
the $X_S$-residualised versions of $X_j$ and $X_i$; and
$\hat r_n^2 = S_{xy}^2 / (S_{xx} S_{yy})$ is the squared sample partial
correlation.

\begin{proposition}\label{prop:safelinear}
Under the linear-Gaussian model~\eqref{eq:regform} with $\beta = 0$, the
process $(E_n)$ of~\eqref{eq:safelinear} is a test martingale with respect to
the filtration generated by the observations, and hence an e-process.
\end{proposition}

The proof (Appendix~\ref{app:safelinear}) is the standard right-Haar argument
\citep{berger1998bayes,hendriksen2021optional}: the null model is invariant
under the group of location shifts in the $X_S$ directions and scale changes,
the right-Haar prior is relatively invariant under that group, and the
resulting Bayes factor is a ratio of marginal likelihoods of the maximal
invariant, whose null distribution is free of the nuisance parameters. Bayes
factors of this form are exactly the ones for which optional stopping is
valid, which is why not every Bayes factor would do here.

Two properties of~\eqref{eq:safelinear} drive the implementation.

\paragraph{Sufficiency.} $E_n$ depends on the data only through Schur
complements of the running Gram matrix
$\sum_{s \le n} (o_s - \bar o)(o_s - \bar o)^\top$. A single $O(d^2)$
sufficient statistic therefore serves the \emph{entire} hypothesis family
$\{(i,j,S)\}$ simultaneously: no per-hypothesis state is stored, and the cost
of an update is $O(d^2)$ regardless of how many hypotheses are live. This is
what makes it feasible to keep every pair's minimum
in~\eqref{eq:adjacency} current after every batch.

\paragraph{The prior scale must be fixed in advance.} $g$ enters
through $G = g S_{yy}$, and $S_{yy}$ grows with $n$. A prior scale chosen to
adapt to the running sample size---for instance by rescaling $g$ so that $G$
stays constant---destroys the martingale property, because the ``prior'' would
then depend on the data it is being used to weigh. We fix $g$ from a warm-up
prefix of the stream that is excluded from every e-process, using it only to
set standardisation constants. The warm-up length appears in the experiments
as a tuning parameter and nowhere in the guarantee.

\begin{remark}[Calibration was verified, not assumed]\label{rem:calibration}
Proposition~\ref{prop:safelinear} is a theorem, but the implementation of a
theorem can be wrong, and a miscalibrated primitive would invalidate every
downstream claim while leaving all the proofs intact. We therefore verified
calibration by simulation: monitoring~\eqref{eq:safelinear} after every
observation over $3000$ replications and $1500$ time points under a true null
gave realised crossing rates of $0.149$, $0.076$, $0.041$ and $0.008$ against
nominal $\alpha \in \{0.20, 0.10, 0.05, 0.01\}$.
Section~\ref{sec:exp-primitive} repeats the exercise under six noise laws.
\end{remark}

\subsection{Confidence sequence by inversion}\label{sec:cs}

Because~\eqref{eq:safelinear} is available in closed form for every value of
the tested coefficient, it can be inverted to give a confidence sequence---a
sequence of intervals covering $\beta$ simultaneously at all times with
probability $1-\alpha$. Writing $\hat\beta = S_{xy}/S_{yy}$ for the
least-squares estimate of the coefficient of $X_j$ in the residualised
regression, the level-$(1-\alpha)$ confidence sequence is
\begin{equation}\label{eq:cs}
\hat\beta \;\pm\; \frac{1}{S_{yy}}
  \sqrt{\frac{\tau\,\bigl(S_{xx} S_{yy} - S_{xy}^2\bigr)}{1-\tau}},
\qquad
\tau \;=\; \frac{1}{\kappa}
  \left[ 1 - \left(\frac{\alpha}{\sqrt{1+G}}\right)^{2/(n-p)} \right] ,
\end{equation}
valid whenever $\tau > 0$, and the whole line otherwise. The derivation is in
Appendix~\ref{app:cs}. We report it because it is the natural summary of
evidence about an edge's strength, because it is what the equivalence-testing
extension of Section~\ref{sec:negligible} uses, and because it provides an
independent check on the implementation: \eqref{eq:cs} must agree with a
brute-force inversion of~\eqref{eq:safelinear}, and an early version of our
code did not. The discrepancy was a sign error in the $\log(1+G)$ term, which
produced intervals roughly half the correct width---a defect invisible in any
Type-I error simulation of the test itself, since the test was correct and
only the reported interval was wrong. Agreement between the closed form and a
numerical inversion is now a unit test.

\subsection{Beyond the linear-Gaussian primitive}\label{sec:beyondlg}

Conditional independence testing is hard in general: \citet{shah2020hardness}
show that no test can have non-trivial power against all alternatives while
controlling size over all null distributions with continuous conditional
laws. Every usable primitive therefore buys power with a model. The two
mainstream routes are the Model-X framework
\citep{candes2018panning}, which assumes the conditional law of $X_i$ given
$X_S$ is known, and parametric or kernel-based assumptions on the joint law.
Sequential conditional-independence tests to date have largely taken the
Model-X route; \citet{csillag2025prediction} note that the requirement is
typically inaccessible in a discovery setting, where the conditional laws are
precisely what is unknown. Construction~\eqref{eq:safelinear} avoids it by
assuming linearity and Gaussianity instead, which is a strong assumption but a
checkable one. Section~\ref{sec:exp-violations} reports what happens when it
fails.
